% !TEX root = ../paper.tex

\section{Results}

\subsection{Pairwise Baseline Comparison}

We begin by asking whether the main difficulty on UAV-TAlign-1K still lies in obtaining usable
pairwise RGB-thermal correspondences. Table~\ref{tab:pairwise_baselines} shows that this is no
longer the dominant bottleneck. Strong off-the-shelf multimodal methods already provide high
pairwise usability, with raw MINIMA and RoMa both reaching 500/500 usable pairwise registrations,
LoFTR reaching 499/500 under a resize-aware inference setting, and XoFTR-640 reaching 494/500.
Classical baselines remain weaker, but they are not uniformly unusable. Taken together, these
results indicate that pairwise multimodal matching availability is already relatively strong on this
dataset.

\begin{table}[t]
\caption{Pairwise baseline comparison on 500 RGB-thermal pairs.}
\label{tab:pairwise_baselines}
\centering
\resizebox{\textwidth}{!}{
\begin{tabular}{lrrrrrrr}
\toprule
Method & OK / N & OK Rate (\%) & Homography (\%) & Mean Inlier Ratio & Mean Coverage & Mean Reproj. Error & Mean Runtime (s) \\
\midrule
raw MINIMA & 500/500 & 100.0 & 100.0 & 0.378 & 0.996 & 1.707 & 1.859 \\
Official pretrained RoMa via local wrapper & 500/500 & 100.0 & 100.0 & 0.400 & 0.982 & 1.672 & 1.077 \\
Kornia LoFTR (pretrained="outdoor") & 499/500 & 99.8 & 99.8 & 0.096 & 0.919 & 1.411 & 0.549 \\
Official pretrained XoFTR-640 via local wrapper & 494/500 & 98.8 & 98.8 & 0.151 & 0.951 & 1.610 & 1.204 \\
AKAZE + RANSAC & 479/500 & 95.8 & 95.8 & 0.177 & 0.646 & 0.393 & 2.368 \\
SIFT + RANSAC & 457/500 & 91.4 & 91.4 & 0.240 & 0.599 & 0.260 & 4.605 \\
\bottomrule
\end{tabular}}
\end{table}

\subsection{Scene-Level Performance of UAV-TAlign}

However, this near-saturated pairwise picture does not carry over to scene-level performance. Under
the canonical QA-aware scene criterion, UAV-TAlign full pipeline passes 7 of 15 scenes, with a
mean accepted-frame ratio of 64.5\%. This gap between strong pairwise usability and substantially
lower scene-level canonical pass rate suggests that the core challenge in UAV RGB-thermal
registration is not simply producing pairwise matches, but forming scene-level transforms that
remain stable after aggregation and reliability-aware acceptance.

\begin{table}[t]
\caption{Formal scene-level result for UAV-TAlign full pipeline.}
\label{tab:scene_level_main}
\centering
\begin{tabular}{lrrrr}
\toprule
Method & Eval Unit & Scene Pass & Pass Rate (\%) & Mean Accepted Ratio (\%) \\
\midrule
UAV-TAlign full pipeline & 15 scenes & 7/15 & 46.7 & 64.5 \\
\bottomrule
\end{tabular}
\end{table}

\subsection{Condition Breakdown}

The scene and condition breakdowns show that the remaining failures are unevenly distributed rather
than uniformly random. Night scenes remain the hardest group under the current canonical gate,
while low-light scenes are mixed rather than uniformly poor. These patterns support the view that
the bottleneck is scene-level reliability rather than pairwise matcher availability alone.

\begin{table}[t]
\caption{Scene-level breakdown by light condition.}
\label{tab:light_condition}
\centering
\begin{tabular}{lrrrr}
\toprule
Condition & Scenes & UAV Pass & Pass Rate (\%) & Mean Accepted Ratio (\%) \\
\midrule
Day & 7 & 4 & 57.1 & 77.1 \\
Night & 5 & 1 & 20.0 & 50.6 \\
Low-light & 3 & 2 & 66.7 & 58.5 \\
\bottomrule
\end{tabular}
\end{table}

\subsection{Protocol and Proxy Evidence}

Because this PRCV version does not include manual geometric ground truth, we examine whether the
stored QA and proxy signals are at least internally consistent with scene-level pass/fail
decisions. The strongest currently available discriminator is the warning code
\texttt{baseline\_relative\_qa\_outliers}, which appears in all 8 failing scenes and in none of the
7 passing scenes. Accepted ratio also correlates with failure, but it is not sufficient by itself:
for example, scene 02 still shows 98\% accepted frames yet fails because reject and QA-relative
outlier signals remain high, whereas scene 04 passes with 80\% accepted frames when those fail-side
warnings are absent. These observations support the claim that the protocol is capturing
scene-level reliability rather than merely counting matched frames.

\subsection{Cumulative Ablation}

The cumulative ablations show that the clearest empirical gains come from robust aggregation and
QA-aware scene decision. Relative to A1, A2 keeps the same accepted-only scene survivability but
substantially improves edge and gradient alignment proxies while reducing severe outliers. A3 then
brings the ablation line into close agreement with the formal subset result, with both A3 and
A4-subset reaching 5/8 canonical scene passes.

\begin{table}[t]
\caption{Main cumulative ablation line on the fixed 8-scene subset.}
\label{tab:ablation_main}
\centering
\resizebox{\textwidth}{!}{
\begin{tabular}{lrrrrrr}
\toprule
Stage & Pass / N & Pass Rate (\%) & Mean Accepted Ratio (\%) & Mean $\Delta$Edge F1 & Mean $\Delta$Grad NCC & Mean Severe-Outlier Ratio (\%) \\
\midrule
A1 & 8/8 & 100.0 & 77.428 & 0.1241 & 0.1261 & 10.417 \\
A2 & 8/8 & 100.0 & 77.428 & 0.1452 & 0.1482 & 4.167 \\
A3 & 5/8 & 62.5 & 77.234 & 0.1465 & 0.1490 & 4.167 \\
A4-subset & 5/8 & 62.5 & 80.387 & 0.1469 & 0.1485 & 4.167 \\
\bottomrule
\end{tabular}}
\end{table}

The completed random-selection supplement further shows a 4/8 to 7/8 spread across seeds on the
same fixed subset, so deterministic-even selection is best written as a reproducible default
operating policy rather than as the single largest performance driver.
